\section{Étape 2 : séparation affichage / calcul avec MPI}

\subsection{Objectif}

L'objectif de cette étape est d'isoler le travail graphique du travail numérique. Au lieu d'exécuter rendu et calcul dans le même processus, on utilise ici exactement deux processus MPI :
\begin{itemize}
    \item le \texttt{rank 0} gère uniquement la fenêtre, les événements SDL et le rendu OpenGL ;
    \item le \texttt{rank 1} conserve tout l'état du système physique et exécute \texttt{update\_positions(dt)}.
\end{itemize}

Cette organisation prépare naturellement la suite du sujet : avant de distribuer réellement le calcul entre plusieurs processus, on vérifie déjà que l'affichage peut être extrait du chemin critique.

\subsection{Changements clés d'implémentation}

\begin{itemize}
    \item un nouveau script MPI a été créé dans \path{stages/02_mpi_display_compute/}, avec exactement deux processus ;
    \item le \texttt{rank 0} ne calcule plus la physique : il reçoit les positions et met uniquement à jour le visualiseur ;
    \item le \texttt{rank 1} conserve tout l'objet \texttt{NBodySystem} et réutilise les noyaux Numba de l'étape 1 pour calculer les pas de temps ;
\end{itemize}

\subsection{Pourquoi exactement deux processus}

Cette étape ne cherche pas encore à paralléliser le calcul par domaine spatial. Le but est plus simple : séparer proprement les responsabilités. Deux processus suffisent donc :
\begin{itemize}
    \item un processus d'affichage ;
    \item un processus de calcul.
\end{itemize}

Utiliser plus de processus ici ajouterait de la complexité MPI sans répondre à l'objectif de l'étape.

\subsection{Protocole de communication}

Le protocole MPI retenu est volontairement minimal.

\begin{itemize}
    \item Au démarrage, le \texttt{rank 1} envoie au \texttt{rank 0} les données nécessaires à l'affichage initial : positions, couleurs, intensités lumineuses et bornes de la scène.
    \item À chaque pas interactif, le \texttt{rank 0} poste d'abord un \texttt{Irecv} pour le prochain tableau de positions, puis envoie la commande \texttt{STEP}.
    \item Le \texttt{rank 1} reçoit \texttt{STEP}, calcule le pas suivant, puis envoie les nouvelles positions.
    \item À la fermeture, le \texttt{rank 0} envoie \texttt{STOP} et attend un accusé de réception.
\end{itemize}

Le choix \texttt{Irecv} puis \texttt{Send} suit directement le schéma vu en cours : le récepteur se déclare prêt avant l'envoi, ce qui évite les interblocages inutiles.